\section{Research Goals}
\label{sec:research_objectives_research_goal}

The research is organised around three goals that build on one another sequentially:
the first characterises the problem, the second identifies its cause, and the third
validates the solution under real-world conditions.

\begin{enumerate}

    \item \textbf{Characterise the Composability Gap.}
    Logical composition is most reliable when the composed conditions affect the data distribution in approximately separable ways.
    When dependencies or interactions between conditions are pronounced, composition can fail: the resulting outputs may not faithfully represent either condition, or may blend them incoherently.
    Our goal is to identify the
    conditions under which the gap between logical and semantic composition is
    significant, and to establish a principled framework for measuring it.
    Closing this gap would enable composable inference to generate complex, unseen
    combinations of conditions without requiring large joint training datasets.

    \item \textbf{Determine the Cause of the Composability Gap and How to Close It.}
    Once the gap is characterised, the central question is what causes it and what is
    required to close it.
    The second goal is to determine whether the gap arises from how concepts are
    represented in latent space, how they are combined during inference, or both.
    This matters because it identifies where intervention is needed and whether
    reliable composition depends on better inference rules, better representations,
    or a combination of the two.
    A successful outcome is a clear account of the mechanism behind the gap and a
    principled basis for deciding how compositional generative models should be
    designed to produce semantically faithful outputs.

    \item \textbf{Demonstrate the Practical Utility of Factorised Composition.}
    To demonstrate that factorised composition can be used to generate novel and
    semantically meaningful outputs that are not easily achievable through large models.
    We restrict our domain to medical imaging, as it presents conditions that are
    well-suited to demonstrating the value of compositional control.
    Disease conditions are discrete and annotatable, and their co-occurrence is
    clinically meaningful.
    At the same time, rare combinations are chronically underrepresented in
    available data.
    The goal is to show that representing disease conditions as independent latent
    concepts enables the generation of ecologically valid co-morbid presentations by
    composing per-pathology models during inference without having seen those
    combinations during training.
    This is particularly valuable in low-resource settings where co-morbid presentations
    are absent from available data, such as silicosis-tuberculosis.
    The broader claim is that approximate clinical independence between conditions is
    sufficient for this to work: exact mathematical independence is not required.

\end{enumerate}
